problem:  
Let \( M^{4n} \) be a compact, oriented Riemannian manifold of dimension \(4n\) with strictly positive sectional curvature.  
Let \( Q_M \colon H^{2n}(M;\mathbb{R}) \times H^{2n}(M;\mathbb{R}) \to \mathbb{R} \) be the intersection form \( Q_M(\alpha,\beta) = \int_M \alpha \wedge \beta. \)  
Denote by \( \sigma(M) \) the signature of this form, i.e. the difference between the numbers of positive and negative eigenvalues of \( Q_M \).

**Conjecture (Hopf–Signature Rigidity Problem):**  
Does positive sectional curvature impose a sign or bound condition on the signature \( \sigma(M) \)? In particular, is it true that
\[
\sigma(M) \ge 0,
\]
and possibly that
\[
|\sigma(M)| \le f(n)
\]
for some universal function \( f \) depending only on the dimension \(4n\)?  
Equivalently: can the algebraic type of the intersection form—especially the presence of negative directions—be limited by the existence of a positively curved metric?

Why is it a "good" problem:  
This conjecture represents a natural, mathematically precise deepening of Hopf’s question about positive curvature and topology. Hopf’s original conjecture restricts the sign of the Euler characteristic in even dimensions—it predicts only that an integrated curvature expression has positive total sign. The present problem asks whether sectional curvature positivity controls the **signature** of the manifold, tying curvature not just to one scalar invariant but to the algebraic structure of the entire middle-dimensional cohomology through the intersection form.

The problem engages multiple fundamental areas:
- **Riemannian geometry and curvature operators:** any attempt to prove such a restriction would involve analyzing how sectional curvature influences the curvature operator acting on \(2n\)-forms, possibly through refined Bochner–Weitzenböck formulas or curvature–operator positivity.
- **Topology and index theory:** the signature is an index invariant, computable via the Hirzebruch L-class and therefore intimately connected with curvature through characteristic classes.
- **Global analysis and higher-dimensional generalization of four-manifold phenomena:** in dimension 4, the sign of the signature corresponds to self-duality properties of the curvature; extending any such relation to higher \(4n\)-dimensions would unify disparate strands of curvature and topology.

All known examples of positively curved \(4n\)-manifolds (such as \( \mathbb{C}P^{2n} \), \( \mathbb{H}P^n \), and certain biquotients) happen to have nonnegative or positive signatures, with intersection forms of definite type. However, no general theorem explains this phenomenon.

The statement is simple, concise, and comprehensible—“Does positive curvature force nonnegative signature or bound it?”—yet its resolution would require a breakthrough combining differential geometry, topology, and index theory. This makes it not only mathematically sound but also a "good" and profoundly forward-looking research problem.